\section{Overlap removal}
\label{sec:ch6_overlap}

\subsection{Shared signals and matching}

The same detector activity can be reconstructed as more than one object.
For example, an electron shower can also form a jet, while a hadronic tau candidate is seeded by a calorimeter jet.
An ordered overlap-removal procedure is applied to the baseline candidates to resolve these ambiguities before event selection.
Only objects surviving an earlier step are considered in subsequent steps.

Geometrical overlaps are evaluated with the rapidity distance $\Delta R_y$ defined in Eq.~\eqref{eq:ch6_rapidity}.
Shared ID tracks are also used to identify duplicate electron and muon interpretations.
For muon--jet matching, \emph{ghost association} is used: an object with negligible momentum is placed along the muon direction during jet clustering, and the jet containing it is identified.
The association is therefore determined by the clustering without changing the jet momentum.

\subsection{Removal sequence}

The following steps are applied in order:
\begin{enumerate}
    \item \textbf{Electron--electron}: when two electrons share an ID track, the electron with lower $\pT$ is removed.
    \item \textbf{Tau--electron}: a tau candidate within $\Delta R_y<0.2$ of an electron is removed.
    \item \textbf{Tau--muon}: a tau candidate within $\Delta R_y<0.2$ of a muon is removed.
    \item \textbf{Calorimeter-tagged muon--electron}: a calorimeter-tagged muon sharing the electron ID track is removed.
    \item \textbf{Electron--muon}: an electron sharing an ID track with a remaining muon is removed.
    \item \textbf{Jet--electron}: a jet within $\Delta R_y<0.2$ of an electron is removed.
    \item \textbf{Electron--jet}: an electron within $\Delta R_y<0.4$ of a remaining jet is removed.
    \item \textbf{Jet--muon}: a jet with fewer than three associated tracks is removed if the muon is ghost-associated with it.
    \item \textbf{Muon--jet}: a muon within $\Delta R_y<0.4$ of a remaining jet is removed.
    \item \textbf{Jet--tau}: a jet within $\Delta R_y<0.2$ of a remaining tau candidate is removed.
\end{enumerate}

The inner electron--jet step retains the electron interpretation of an electromagnetic shower, whereas the outer step rejects electrons embedded in surviving hadronic jets.
For muon--jet overlaps, the small track multiplicity is used to recognise a jet candidate dominated by activity associated with the muon.
The last step retains the tau interpretation of an overlapping jet.
The separate low-$\pT$ muon rejection in tau identification, described in Section~\ref{sec:ch6_tau}, is applied in addition to this baseline-object sequence.
